% !TeX root = ../papers/nc-p3.tex
\section{片側対称因子を横切る中心の二次・四次完全交叉}
\label{sec:symmetric-center-cover}
以下は偶二位置が対称となる因子$D$を扱う。基底付きテンソルの滑らかな成分を
$r$の近くで考え、中心の次数は経路長の四分の一で表示する。
$Z_j$をこの換算での次数$j$部分とする。

\begin{theorem}\label{thm:symmetric-center-cover}
$r$を含む成分の開近傍の全幾何点は標準Koszul左右AS正則である。
基底を局所的に選べば、$P=k[x_1,\ldots,x_4]$上の全中心は
$1,z,y,zy$を自由基底に持ち、$\deg z=1$, $\deg y=2$である。その表示は
\[
 z^2=a_2+a_1z+\tau y,\qquad
 y^2=b_4+b_3z+b_2y+b_1zy,\qquad a_i,b_i\in P_i .
\]
$\tau$は$D$の局所方程式である。因子上では$b_1=b_2=0$となるように
$y$を選べる。従って$D$上の射影中心は滑らかな三次元二次超曲面の二重被覆で、
分岐は四次式で切られ、経路の反転が被覆変換を誘導する。因子外では$y$を
消去して四次超曲面を得る。近傍の全幾何点で中心は幾何学的に整かつ正規、
折り畳み環のPI次数は16で、各隅の商可除代数の次数は4である。
\end{theorem}
\begin{proof}
$r$の四次中心の次元は5で、うち四切断は既出の多項式部分環を与える。
十五次元成分の一般点でも次元5であるから、経路長4,5の標準階数の開集合で
可換方程式の核は$r$の近くでも階数5の部分束である。四変数で割った次数13までの
判定をその切断に課せば、定理\ref{thm:full-component-central-open}の自由性と
左右完全性の証明がそのまま$r$を含む近傍に適用される。
一般点の四次超曲面表示から$\dim Z_2=15$を知るので、定階数の経路空間上の
可換方程式により$r$では少なくとも15である。

固定同定上、反転は長さ$4j$の語を逆順に読み、始点$i$を$3-i\pmod4$に送る。
二次関係の四つの表示からこれは反自己同型であり、中心では自己同型となる。
対称軌跡での同定の局所回復は定理\ref{thm:full-fifteen-component}による。
$K$を対角四元群の文字群とする。$r$で各矢と可換な頂点対角元を、
文字別に経路長9までの商基底で計算する。素数1009における
反転の正負固有空間の次元は
\[
\begin{array}{c|cccc}
 &0&1&2&3\\\hline
 Z_1&(3,0)&(2,0)&(0,0)&(0,0)\\
 Z_2&(8,1)&(6,0)&(0,0)&(0,0)
\end{array}
\]
である。商基底は右端から一文字ずつ延長して六関係の右倍を掃き出す。
各次数の経路次元は検算範囲で標準値と一致する
（各頂点で経路長$n$の次元は$\binom{n+3}{3}$）。
可換核の上界の合計が既知の下界と一致するため、局所整数環でも核は
直和因子で剰余への写像は全射である。$2$が可逆なので固有空間の射影も
特殊化と可換である。従って表は標数零で正しく、近傍の$Z_2$は階数15の束となる。
$D$上で反転不変な$Z_1$の四変数と第五切断$z$を取り、$Z_2$の
唯一の反不変直線の局所基底$y$を取る。この固有次元は因子上で局所一定である。
これらの切断を成分の近傍へ延ばす。

前節のフィルター付き有限商の基底と安定した矢作用で$z,y$を作用させる。
四変数の値を零とした第0隅で$z,y,zy$の乗法階数は剰余で順に32,2,1である。
特に次数の異なる$1,z,y,zy$の像は独立で、$P$自由な折り畳み環の
自由直和因子を張る。四変数の値$(1,2,3,4)$では剰余で可換核次元4、
第0隅の正則trace対合の階数64を得る。前節と同じ整数環からの階数比較と
分数体上の飽和性により、この四元が全中心の$P$自由基底になる。
斉次性により定理の二式を得る。これらが全関係を生成することは、
$P$上で$z,y$に比較用の重み2,3を付ければ従う。左辺の平方を右辺で
書き換えると重みが下がり、残る四単項式は独立である。
$D$上の固有符号から$\tau=b_1=b_2=0$が従う。

次に$\mu:\operatorname{Sym}^2 Z_1\to Z_2$を考える。
四変数の二次単項式十個と$x_i z$の四個は自由基底表示で独立である。
残る一列$z^2$の余核方向への係数が$\tau$なので、
$\det\mu$は単元倍の$\tau$である。$r$の唯一の横断商接方向で
その微分が非零であることを確認する。節\ref{sec:unrestricted-doublet}で
有理数上の軌道方向を除いて取った文字0の交代代表を$T$とする。
$p=1009$とし、$W+pT,W-pT,W^{\mathsf T}+pT^{\mathsf T},
W^{\mathsf T}-pT^{\mathsf T}$の六基底列を$\bmod\ p^2$で使う
（奇位置は二添字を交換）。経路長8までの商と長さ4の可換核は単元ピボットで自由である。
この核の文字0の三元と文字1の二元の全積を商に投影する。
文字0の九列には八つの単元ピボットがあり、残部は$p$で割れて$p^2$で
消えない。文字1の六列には六ピボットがある。
基点で階数14であることは既に標数零で分かった。これらの基底と掃き出しは
局所整数環上で単元のみを割る操作であり、基点で消える最大小行列の
$W+pT$での値は$p$とその微分の積に法$p^2$で等しい。
従って像を経路商に入れても非零の法線微分があり、$d\tau(T)\ne0$である。
滑らかな因子$D$上で$\tau=0$であったから、$\tau$はその局所方程式の単元倍となる。

$D$上では二次超曲面の座標環は$P\oplus Pz$、全中心はその上で$1,y$を
基底とする。分岐式は$g=b_4+b_3z$である。四変数を
$(s,t,s+2t,3s-t)$に制限する。$s=1$で$t=0,1,2,3,4$の五つの有限商の
正則作用から$z^2,y^2$を$1,z$で表示して補間すると、$r$の剰余では平面の
二次式のHesse行列が階数3となる。二次式と$g$に偏微分の二行二列小行列を
加え、全射影座標層で共通零点がないことを検算する。
従って標数零でも滑らかな平面二次曲線上で八個の分岐点は単純であり、
その被覆は滑らかで幾何学的に整な種数3の曲線である。
二次超曲面自身についても、$(1,i,j,k)$, $i+j+k\le2$の十点（非負整数）で
$z^2$の係数を補間すると五変数のHesse行列の階数は5であり、基底二次超曲面は滑らかである。
正確な行列計算は\path{computations/verify_symmetric_center.py}にある。

同じ二線型切断上の曲線の滑らかさと幾何学的整性は近傍で開く。
実際、二式で定まる重み付き射影族は四変数の射影空間上有限平坦であり、
指定した直線上への制限も有限平坦かつ固有である。重み2の頂点を通らないので
通常のJacobi判定を使える。全射影中心の各成分は三次元で、
特異軌跡が二次元以上なら有限射影で指定直線に会い、滑らかな曲線切断に反する。
アフィン錐も頂点以外は同じ正則性判定が通る。完全交叉ゆえに$S_2$、
従って正規で、次数0部分が体の正次数付き環には非自明な冪等元がないため整となる。
前節のtraceと斉次右イデアルの階数による議論が近傍全体で適用できる。
\end{proof}
